\documentclass[11pt,a4paper]{article}
\usepackage[margin=25mm]{geometry}
\usepackage[T1]{fontenc}
\usepackage[utf8]{inputenc}
\usepackage{lmodern}
\usepackage{amsmath,amssymb}
\usepackage{booktabs}
\usepackage{array}
\usepackage{longtable}
\usepackage{microtype}
\usepackage[hidelinks]{hyperref}
\usepackage{url}
\setlength{\emergencystretch}{2em}
\newcommand{\Expect}{\mathbb{E}}
\newcommand{\artifact}[1]{\path{#1}}
\newcommand{\R}{\mathcal{R}}

\title{A Diffusion-Guided Neural Architecture for Capacitated Vehicle Routing\\
with Classical and QUBO-Based Local Refinement}
\author{Merve, Hassan and Mustafa Mert}
\date{7 October 2026}

\begin{document}
\maketitle

\begin{abstract}
We present an implemented neural solver for the capacitated vehicle routing
problem, inspired by Constraints Matrix Diffusion based Generative Neural
Solver for Vehicle Routing Problems. Our architecture combines a discrete
diffusion model of customer route membership with a global Transformer,
a matrix-masked local Transformer, learned gated fusion, and a capacity-aware
dual-pointer decoder trained by multi-start reinforcement learning. The
predicted matrix supplies structural guidance, while explicit decoding masks
enforce feasibility. A modular improvement stage reuses the matrix for
neighborhood selection and maps its probabilities into local quadratic
unconstrained binary optimization (QUBO) objectives. We evaluate the frozen
solver and classical, annealing, and simulated quantum refinement on 96
reserved instances and 72 spatial stress instances with 20, 50, or 100
customers. On the reserved set, the constructive solver has a mean gap of
19.22\% to strengthened heuristic references. Inexpensive classical refinement
reduces this to 9.69\%, and larger restarted neighborhoods reach 4.14\%;
the spatial stress gap falls from 16.53\% to 3.54\%. Matrix-guided selection
identifies useful exchanges, and matrix-derived objective terms improve
annealing in exploratory ablations. Simulated QAOA with at most 16 qubits gives
no statistically resolved improvement over exact classical solving or matched
random sampling. The study establishes a complete construction-and-refinement
architecture and measures its current limitations; it does not demonstrate
quantum advantage or state-of-the-art CVRP performance.
\end{abstract}

\noindent\textbf{Keywords:} capacitated vehicle routing; discrete diffusion;
neural combinatorial optimization; gated fusion; reinforcement learning;
local search; QUBO.

\section{Introduction}

Capacitated vehicle routing couples two decisions: which customers share a
vehicle and in what order that vehicle visits them. A constructive neural
solver makes these decisions sequentially; an improvement method revisits
them after a feasible solution is available. Learned information about
customer grouping can potentially help at both stages, provided that it is
combined with an explicit mechanism for enforcing routing constraints.

This article presents our implemented construction-and-refinement solver.
Its inspiration is the constraints matrix diffusion (CMD) approach of
Wang et al.~\cite{wang2026cmd}, which generates a structural assignment
matrix and integrates it into an autoregressive routing network. We adopt
the idea of predicting same-route structure, then build a solver around
a trainable global Transformer, a matrix-masked local Transformer, gated
feature fusion, and a dual-pointer decoder. A multi-start training objective
and explicit visited-customer and capacity masks complete the constructive
architecture.

We also give the predicted matrix a role after construction. It ranks
neighborhoods that may contain route-membership errors and supplies
coefficients for local QUBO objectives. The refinement stage is modular:
the same starting routes can be improved by classical search, simulated
annealing, simulated quantum annealing, or simulated QAOA. This makes the
architecture a controlled test bed for learning-guided local optimization,
even when classical methods remain the strongest solvers.

\subsection{Contributions and scope}

The work has three contributions:
\begin{enumerate}
 \item \textbf{An implemented diffusion-guided CVRP architecture.}
 We describe the complete path from audited route-membership supervision
 to diffusion inference, global and local representation learning, gated
 fusion, and feasible multi-start route construction. The architecture is
 our implementation inspired by CMD; its standard attention, diffusion,
 and reinforcement-learning components are credited to their sources.
 \item \textbf{A reusable local refinement interface.}
 The prior guides neighborhood selection and can be converted into linear
 assignment costs and quadratic customer-pair interactions. Reorder and
 exchange operators share extraction, reconstruction, repair, and true-cost
 acceptance across classical and QUBO solvers.
 \item \textbf{An empirical evaluation of the complete solver.}
 Frozen checkpoints, heuristic-reference auditing, reserved and spatial
 stress panels, fixed-subproblem comparisons, matrix ablations, and
 matched random controls quantify construction quality, refinement value,
 and the limits of simulated quantum optimization.
\end{enumerate}

Three questions organize the evaluation: how well does the implemented
constructive architecture solve CVRP, how much does local refinement improve
its solutions, and does matrix-guided QUBO optimization add value relative
to classical controls? We report the frozen implementation rather than
attributing every gain to an individual architectural component. The
available experiments establish pipeline behavior and selected guidance
effects; they do not isolate the causal contribution of every encoder,
gate, or decoder feature.

\section{Related work}
\label{sec:related}

\subsection{Neural routing and structural diffusion}

Attention-based routing policies learn distributions over sequential
customer visits. Kool et al.~\cite{kool2019attention} train attention
models with REINFORCE for CVRP and other routing problems. POMO
\cite{kwon2020pomo} uses diverse starts and a shared baseline to exploit
alternative solution representations. Our constructive policy uses a
multi-start objective and keeps the best greedy start at inference.

Discrete diffusion learns denoising transitions in a finite state space
\cite{austin2021discrete}. DIFUSCO~\cite{sun2023difusco} applies graph-based
diffusion to binary representations of traveling-salesperson and
independent-set problems. CMD~\cite{wang2026cmd} instead generates a
constraints assignment matrix for routing and integrates it into a fusion
network and autoregressive solver. That idea motivates our membership
prior. Our implementation uses Transformer-based global and local policy
encoders with a learned gate and extends the structural prior into
post-construction refinement.

Membership and adjacency are different prediction tasks. A high
same-route probability supports assigning two customers to one vehicle;
it does not imply that they should be consecutive visits. This distinction
informs both our policy masks and the interpretation of matrix-derived
reorder costs.

\subsection{Classical and QUBO-based optimization}

Hybrid genetic search combines population search with strong local
improvement. Vidal's CVRP implementation and SWAP* neighborhood illustrate
the importance of effective customer exchanges~\cite{vidal2022hgs}.
PyVRP provides a high-performance hybrid genetic search package
\cite{wouda2024pyvrp} and supplies our heuristic references. Our local
controls use simpler 2-opt, relocate/swap, restarts, and exact solving.
For short segments, Held--Karp dynamic programming~\cite{held1962dynamic}
is an especially relevant exact control because it searches valid orders
without introducing one-hot infeasible states.

Quadratic binary and Ising formulations encode constraints through penalty
terms~\cite{lucas2014ising}. QAOA alternates objective and mixing unitaries
with classically optimized parameters~\cite{farhi2014qaoa}.
Azad et al.~\cite{azad2020vrp} study Ising formulations and simulated QAOA
for small routing problems. Here, QUBOs encode bounded neighborhoods of
an existing solution. Learned matrix entries change their coefficients,
and common classical postprocessing maps samples back into routes.
This permits controlled tests against exact search and uniform sampling.
Our contribution is the implemented architecture, its refinement interface,
and their evaluation; we do not claim priority for neural routing,
QUBO routing, or hybrid quantum optimization in general.

\section{Problem definition and structural representation}
\label{sec:problem}

Let the depot be node $0$, let customers be $1,\ldots,n$, and let
$v_i\in\mathbb R^2$ and $q_i$ be their coordinates and demands. Distances
are $d_{ij}=\|v_i-v_j\|_2$. A solution $\R$ is a collection of routes
that starts and ends at the depot, visits every customer exactly once,
and satisfies $\sum_{i\in r}q_i\leq Q$ for every route $r$. We minimize
\begin{equation}
 C(\R)=\sum_{r\in\R}\sum_{(i,j)\in r}d_{ij}.
 \label{eq:cvrp-cost}
\end{equation}
The fleet is homogeneous, split deliveries are disallowed, and the number
of used routes may change during refinement. The synthetic experiments
minimize distance without a separate vehicle-use penalty or prescribed
fleet size. They cover $n\in\{20,50,100\}$ and capacities $Q=30,40,50$,
respectively. The random, clustered, and mixed spatial regimes are plain
CVRP; they do not include time windows.

For distinct customers, define $M_{ij}=1$ when $i$ and $j$ belong to the
same reference route and zero otherwise. We set $M_{ii}=0$ by convention
and exclude diagonal entries from the loss. This representation is invariant
to route labels and visit order within a route. It captures customer
partitioning, but does not identify a route's first customer, its internal
sequence, or its depot connections. A soft prediction $P=(p_{ij})$ also
need not form a transitive partition. It is therefore used as a structural
prior rather than treated as a feasible CVRP solution.

Different near-equal-cost solutions can partition customers differently.
Our label audit checks both reference cost and partition stability, and
retains competitive alternatives for stochastic supervision. The three
training pools each contain 500 source instances. Their materialized
partition counts are 500, 536, and 1,268 at N20, N50, and N100,
respectively. Each diffusion training epoch chooses one retained
competitive partition per source, using a
source- and epoch-dependent deterministic selection seed. Alternatives
are thus not counted as independent graphs. Reference costs are heuristic
solutions, not certified optima.

\section{Implemented solver architecture}
\label{sec:architecture}

Our solver has a structural prediction stage, a constructive policy, and
an optional improvement stage. Figure~\ref{fig:pipeline} summarizes the
interfaces. The diffusion stage works on customer pairs. The policy encodes
the depot and all customers, constructs full routes, and enforces feasibility.
Refinement reuses the predicted matrix and improves the resulting incumbent
without changing the neural weights.

\begin{figure}[tbp]
\centering
\fbox{\parbox{0.93\linewidth}{\centering\small
Instance $G$ $\longrightarrow$ graph denoiser $\longrightarrow$
route-membership prior $P$\\[5pt]
$G$ $\longrightarrow$ global Transformer $\longrightarrow$
$P$-masked local Transformer\\[3pt]
$\longrightarrow$ gated fusion $\longrightarrow$ dual-pointer decoder
$\longrightarrow$ best of 16 feasible starts\\[5pt]
$(\R_0,P)$ $\longrightarrow$ selected neighborhoods $\longrightarrow$
classical or QUBO solver\\[3pt]
$\longrightarrow$ reconstruction and feasibility checks
$\longrightarrow$ accepted lower-cost routes
}}
\caption{The implemented architecture. Global embeddings feed the local
encoder and also reach fusion directly. The same fixed prior guides
construction and refinement; feasibility is checked independently of it.}
\label{fig:pipeline}
\end{figure}

\subsection{Discrete diffusion of route membership}

Let $G$ denote instance features and $M^{(0)}$ a clean membership matrix.
The forward process corrupts upper-triangular customer pairs and mirrors
them to preserve symmetry. With $T=700$ transitions, the frozen recipe
uses a linear schedule from $\beta_1=10^{-4}$ to $\beta_T=0.02$ and
\begin{equation}
 Q_t=\begin{pmatrix}1-\beta_t&\beta_t\\
                    \beta_t&1-\beta_t\end{pmatrix},
 \qquad \overline Q_t=Q_1Q_2\cdots Q_t.
 \label{eq:forward-diffusion}
\end{equation}
The marginal flip probability at time $t$ is
$\tfrac12[1-\prod_{u=1}^{t}(1-2\beta_u)]$. Diagonal and padded entries
remain zero. The depot has no membership-matrix row.

The denoiser uses a pretrained frozen graph attention node encoder and
an eight-layer anisotropic graph network. Noisy pair values and geometric
distances initialize the edge features. Sinusoidal timestep embeddings
condition the message-passing layers, and the output predicts clean-pair
logits $f_{\theta,ij}(G,M^{(t)},t)$. The frozen implementation uses width
192, five node-GAT layers with eight heads, and LayerNorm in the denoising
network. Define $\widehat p_{ij}^{(t)}=\operatorname{sigmoid}(f_{\theta,ij})$.
For ordered off-diagonal real pairs $\mathcal V$ in a batch, training uses
\begin{align}
 \mathcal L_{\mathrm{diff}}
 &=\Expect_{G,M^{(0)},t,M^{(t)}}\left[
   \frac{1}{|\mathcal V|}\sum_{(i,j)\in\mathcal V}
   \ell(M_{ij}^{(0)},\widehat p_{ij}^{(t)})\right],\\
 \ell(m,p)&=-w_+m\log p-(1-m)\log(1-p),
 \qquad w_+=\left(\frac{N_-}{\max(N_+,1)}\right)^{1/2}.
 \label{eq:diffusion-loss}
\end{align}
Here $N_+$ and $N_-$ are the batch's positive and negative pair counts,
and $t$ is uniform over training timesteps. This weighted classification
loss predicts clean membership at different corruption levels; it does
not directly minimize route cost.

For a skipped reverse transition from $t$ to $s<t$, inference uses the
learned clean-state mixture
\begin{equation}
 p_\theta(M_{ij}^{(s)}\mid M_{ij}^{(t)},G)=
 \sum_{b\in\{0,1\}}
 q(M_{ij}^{(s)}\mid M_{ij}^{(t)},M_{ij}^{(0)}=b)
 p_\theta(M_{ij}^{(0)}=b\mid M^{(t)},G).
 \label{eq:reverse-mixture}
\end{equation}
The conditional $q$ follows by Bayes' rule from the cumulative forward
kernels. Inference starts from a symmetric Bernoulli matrix with positive
probability $1/2$ and uses 50 reverse steps. The final clean-pair
probabilities form $P$, which is cached throughout refinement. The
implementation names this sampler \texttt{posterior\_mixture\_v2}.
Independent denoiser checkpoints are used for the three customer counts.

\subsection{Global instance encoder}

The global encoder processes the complete depot-and-customer instance.
Its input feature for node $i$ is
\begin{equation}
 u_i=[v_{i,x},v_{i,y},q_i/Q,\mathbf 1\{i=0\}].
 \label{eq:node-features}
\end{equation}
A learned projection maps $u_i$ to width 128. Five Transformer layers
with eight heads and feed-forward width 512 produce global embeddings
$h_i^g$. Only padding is masked in this branch; the predicted membership
matrix does not restrict global attention. This gives the decoder a
complete view of geometry, normalized demand, and depot identity even
when the structural prior is inaccurate.

\subsection{Matrix-masked local encoder}

The local Transformer takes $h_i^g$ as its input and applies five further
attention layers. For distinct customers its frozen attention mask allows
$(i,j)$ if the symmetrized prior satisfies $p_{ij}\geq\tau$, with
$\tau=0.5$. Every real node retains a self-loop. Customer queries may
attend to the depot; the depot query attends only to itself in the local
branch. This asymmetric depot rule supplies depot context without making
the depot a hub for messages between otherwise separate local groups.
The global branch already carries full-instance information.

For one local attention head, the operation can be written as
\begin{equation}
 A^l_{ij}=\frac{(W_Q h_i)^\top(W_K h_j)}{\sqrt{d_h}}
             +\mu_{ij},\qquad
 \mu_{ij}=\begin{cases}0,&(i,j)\text{ is allowed},\\
                     -\infty,&\text{otherwise},\end{cases}
 \label{eq:local-attention}
\end{equation}
followed by row-wise softmax, value aggregation, and residual attention
and feed-forward updates. The output is $h_i^l$. Hard masking concentrates
local messages on predicted route groups; it does not prevent the final
policy from visiting a customer outside the current group.

\subsection{Learned gated fusion}

Rather than replacing the global representation, the implementation
learns a feature-wise gate between aligned global and local embeddings:
\begin{align}
 g_i&=\operatorname{sigmoid}\!\left(
       \operatorname{MLP}([h_i^g,h_i^l,|h_i^l-h_i^g|])\right),\\
 \widetilde h_i&=h_i^g+g_i\odot(h_i^l-h_i^g).
 \label{eq:fusion}
\end{align}
A position-wise residual feed-forward block and normalization produce
fused embeddings $h_i$. Their masked mean, including the depot, gives
the graph context $\overline h$. The gate is initialized to favor the
global branch, with local weight 0.25. It can learn different local
weights for different nodes and features. This is an architectural
mechanism for moderating the prior's influence, not empirical proof that
the trained gate always rejects incorrect matrix information.

\subsection{Capacity-aware dual-pointer decoding}

At decoding step $t$, let $i_t$ be the current node and $L_t$ the remaining
vehicle capacity. The local query context combines depot, current-node,
and projected capacity embeddings; the global query additionally includes
$\overline h$. A global multi-head pointer attends to all feasible nodes.
A local pointer attends to feasible neighbors of $i_t$ in the prior mask,
with depot access retained so that routes can close. Their glimpses are
summed with a learned context-perception term, projected to a query $a_t$,
and used to score candidate embeddings.

The geometric savings term for two customer nodes is
\begin{equation}
 s_{ij}=d_{0i}+d_{0j}-d_{ij}.
 \label{eq:savings}
\end{equation}
With score keys $k_j$, the decoder's customer logits have the form
\begin{equation}
 z_{tj}=10\tanh\!\left(
       \frac{a_t^\top k_j}{\sqrt{128}}
       +\log\max(s_{i_tj},0.01)\right).
 \label{eq:decoder-logit}
\end{equation}
Depot-related savings use a learned scalar bias. Visited customers and
customers with $q_j>L_t$ receive $-\infty$ logits after clipping.
A depot visit closes the current route and resets capacity. The global
pointer keeps feasible customers available even when local membership
predictions are wrong. Thus the structural mask influences preferences,
while visited and capacity masks determine feasible actions.

For instances with $q_i\leq Q$, the policy constructs complete feasible
routes and returns to the depot. The frozen inference protocol forces
16 distinct first-customer starts, selected in increasing depot distance,
decodes greedily thereafter, and keeps the cheapest solution. Encodings
are reused across starts. The feasible incumbent is the interface to
all subsequent refinement methods.

\subsection{Reinforcement learning and training configuration}

The diffusion encoder is trained first, and its prior is held fixed while
the policy learns route construction. Policy training uses a POMO-style
shared baseline~\cite{kwon2020pomo}. For $S$ sampled starts on an instance,
let $C_s$ be rollout cost and $\overline C=S^{-1}\sum_s C_s$. The loss is
\begin{equation}
 \mathcal L_{\mathrm{policy}}=
 \frac1S\sum_{s=1}^{S}\operatorname{stopgrad}(C_s-\overline C)
       \log\pi_\psi(\R_s\mid G,P).
 \label{eq:policy-loss}
\end{equation}
The sign is appropriate for cost minimization: above-average-cost
rollouts have their probabilities reduced. Diffusion training uses Adam;
policy training uses Adam with weight decay $10^{-6}$. The selected
weights are frozen for evaluation and refinement.

Table~\ref{tab:training} gives the archived configurations. Epoch and
runtime limits are configured ceilings, not assertions that every run
completed every epoch. Denoisers use training seed 4331 and policies
seed 42. Diffusion checkpoint selection uses decoded validation route
gap; policy selection uses validation best-start cost. The N100 policy
is the primary seed of the 16-start training experiment, rather than
the seed with the best development score.

\begin{table}[tbp]
\centering\small
\setlength{\tabcolsep}{4pt}
\caption{Frozen training configurations. Denoisers share a recipe but
have separate weights. Global and local policy encoders each have five
layers and eight heads.}
\label{tab:training}
\begin{tabular}{@{}lrrrr@{}}
\toprule
Setting & Diffusion & Policy N20 & Policy N50 & Policy N100\\
\midrule
Hidden/embedding width & 192 & 128 & 128 & 128\\
Denoising layers & 8 & -- & -- & --\\
Node-GAT layers/heads & 5/8 & -- & -- & --\\
Learning rate & $3\cdot10^{-4}$ & $10^{-4}$ & $10^{-4}$ & $10^{-4}$\\
Batch size & 4 & 16 & 16 & 16\\
Gradient accumulation & 4 & 1 & 1 & 1\\
Maximum epochs & 15 & 15 & 8 & 6\\
Training starts & -- & 8 & 8 & 16\\
Time ceiling (s) & 3,600 & 1,800 & 2,400 & 10,800\\
Gradient-norm clipping & 1 & 1 & 1 & 1\\
\bottomrule
\end{tabular}
\end{table}

Dense pair prediction and attention require quadratic storage in customer
count, and the reverse chain evaluates the denoiser repeatedly. Multi-start
decoding then adds sequential policy work. These costs motivate measuring
both constructive inference and optional refinement time, rather than
reporting decoder time alone.

\section{Local refinement extension}
\label{sec:refinement}

\subsection{Neighborhoods and acceptance}

Four selectors identify expensive route segments, customers with uncertain
membership between two routes, consecutive customers with low predicted
same-route probability, and geometrically nearby route pairs. Their proposals
create either a segment-reordering problem or a two-route exchange problem.
All other customers and the relevant segment boundaries remain fixed.
For exchanges, fixed customers retain their route membership, while the
reconstruction may change their visit order within that route.

The uncertainty selector compares a customer's mean predicted affinity
to its own route with its best competing route. If the difference is
below 0.2, the customer is a candidate for reassignment; candidates are
grouped by the implicated route pair. Low-confidence edges use a 0.5
threshold on consecutive-customer probabilities. The geometric selector
considers nearby route pairs and boundary customers, while the high-cost
selector ranks routes by cost per customer and proposes reorder segments.

Every proposed solution is decoded, repaired if necessary, and rescored using
the true full-route cost. The update is accepted only when
\begin{equation}
 \R'\text{ is feasible}\quad\text{and}\quad C(\R')<C(\R).
 \label{eq:accept}
\end{equation}
Thus, an unsuccessful or infeasible QUBO proposal leaves the current solution
unchanged. The loop repeats selection and refinement until a round makes no
improvement or reaches its round limit. This is local refinement of an existing
CVRP solution; the entire instance is never encoded as one QUBO.

\begin{figure}[tbp]
\centering
\fbox{\parbox{0.92\linewidth}{\small
\textbf{Refinement loop.} Input: instance $G$, feasible starting routes
$\R_0$, fixed prior $P$, selector limits, and solver.\\[3pt]
1. Set the incumbent to $\R_0$.\\
2. At each round, rank neighborhoods using the current routes and $P$.\\
3. For each proposal, extract its current subproblem. Run the chosen
solver directly, or build $E_\alpha$, sample binary states, and decode
and repair the retained states.\\
4. Reconstruct affected routes and choose the feasible candidate with
lowest true cost. Replace the incumbent only if Equation~\eqref{eq:accept}
holds; record attempts, repairs, accepted moves, and time.\\
5. Stop after a round with no accepted move or after 30 rounds. Return
the incumbent.}}
\caption{Common improvement loop. Proposals are re-extracted from the
current incumbent before solving so that earlier accepted moves are
respected.}
\label{fig:algorithm}
\end{figure}

Strict cost decrease prevents cycling between incumbents. However,
stopping is relative to the selected neighborhoods and solver budget,
and does not certify CVRP optimality.

\subsection{Route reordering}

For $k$ customers between fixed endpoints $s$ and $e$, let $x_{ip}=1$ when
customer $i$ occupies position $p$. The one-hot formulation uses $k^2$ binary
variables:
\begin{align}
 E_{\mathrm{reorder}}(x)
 &=\sum_i d_{si}x_{i1}
   +\sum_{p=1}^{k-1}\sum_{i\ne j}d_{ij}x_{ip}x_{j,p+1}
   +\sum_i d_{ie}x_{ik} \nonumber\\
 &\quad+A\sum_i\left(\sum_p x_{ip}-1\right)^2
       +A\sum_p\left(\sum_i x_{ip}-1\right)^2.
 \label{eq:reorder}
\end{align}
The default penalty is $A=2\max(d_{uv})$ over the local distance matrix,
with a positive numerical floor, and is enlarged when
bias terms are enabled. On valid permutations, the unbiased energy equals
the segment's true path cost. Invalid samples are repaired by a subset dynamic
program that maximizes agreement with the sampled bits and then minimizes
path cost among equally agreeing permutations. Consequently, classical repair
can contribute substantially to the quality of a sampled solution.
The encoding scales quadratically: four customers require 16 variables
and five require 25. Exact path dynamic programming instead requires
$O(k^2 2^k)$ time and $O(k2^k)$ storage.

\subsection{Two-route exchange}

Let $y_i=1$ assign movable customer $i$ to route $a$ and $y_i=0$ assign it to
route $b$. Let $I_a(i)$ and $I_b(i)$ be cheapest insertion costs into the
fixed customers of those routes. The implemented surrogate is
\begin{equation}
 S(y)=\sum_i\bigl[y_iI_a(i)+(1-y_i)I_b(i)\bigr]
 +\frac{\gamma}{\max(m-1,1)}\sum_{i<j}d_{ij}
       (1-y_i-y_j+2y_iy_j),
 \label{eq:exchange}
\end{equation}
where $m$ is the number of movable customers and $\gamma$ controls dispersion.
The last factor is one when the customers are assigned to the same route.
For fixed loads $F_a,F_b$, capacity penalties have the form
\begin{equation}
 A\left(F_a+\sum_i q_i y_i+u_a-Q\right)^2
 +A\left(F_b+\sum_i q_i(1-y_i)+u_b-Q\right)^2,
\end{equation}
with nonnegative binary-encoded slack $u_a,u_b$. Constraints that cannot bind
need no slack. Integral demands permit exact load representation, and the
penalty is chosen to dominate the surrogate's range.
For an integer slack upper bound $U>0$, the implementation uses
$b=\lfloor\log_2 U\rfloor+1$ bits with coefficients
$1,2,\ldots,2^{b-2},U-(2^{b-1}-1)$; for $b=1$ the coefficient is
just $U$. Their subset sums cover exactly $0,\ldots,U$.
An exchange therefore requires $m+b_a+b_b$ variables, where an inactive
capacity constraint contributes zero bits. Movable-customer count alone
does not determine whether an exchange fits a qubit limit.

With $\gamma=1$, a bound on the unbiased surrogate range is
\begin{equation}
 D_S=\sum_i|I_a(i)-I_b(i)|
 +\frac{\gamma}{\max(m-1,1)}\sum_{i<j}d_{ij}.
 \label{eq:exchange-span}
\end{equation}
The default exchange penalty is $A=1.5\max(D_S+D_B,10^{-9})$, where
$D_B$ bounds the range of the weighted matrix bias. Integral loads imply
that an infeasible assignment has a squared capacity residual of at
least one. The penalty then places such an assignment above feasible
assignments with correctly encoded slack. This ground-state property
does not ensure that finite annealing or QAOA samples are feasible.

Unlike the reorder energy, $S(y)$ is not the true route objective. Decoding
reconstructs each route using nearest-neighbor ordering and 2-opt before
applying Equation~\eqref{eq:accept}. An ``exact exchange'' result therefore
means the best assignment under this reconstruction rule, not an exact joint
solution of membership and visit order. Capacity-violating samples are
repaired by moving a movable customer out of the overloaded route,
choosing a move that fits the receiving route and minimizes the surrogate.
Unrepairable assignments are discarded. Empty routes may be removed,
allowing exchanges to reduce route fragmentation.

\subsection{Proposed diffusion-matrix-to-QUBO mapping}
\label{sec:matrix-qubo}

The proposed formulation augments the preceding geometric objectives with
coefficients computed from the diffusion output $P$. For binary decision
vector $z$, write
\begin{equation}
 E_\alpha(z;P)=C_{\mathrm{local}}(z)+\Pi_\alpha(z)+\alpha B(z;P),
 \qquad \alpha\geq 0,
 \label{eq:diffusion-qubo}
\end{equation}
where $C_{\mathrm{local}}$ is the reorder path cost or exchange surrogate,
$\Pi_\alpha$ contains the constraint penalties, and $B$ expresses agreement
with the learned matrix. The matrix is fixed during local optimization, so
the resulting objective remains quadratic in the decision variables.

For exchanges, let $\mathcal{F}_r$ contain the fixed customers of route
$r\in\{a,b\}$. Define the predicted affinity of movable customer $i$ to
that route by
\begin{equation}
 a_{ir}=\begin{cases}
 |\mathcal{F}_r|^{-1}\displaystyle\sum_{j\in\mathcal{F}_r}p_{ij},
       &\mathcal{F}_r\ne\varnothing,\\
 1/2, &\mathcal{F}_r=\varnothing.
 \end{cases}
\end{equation}
In probability mode the matrix-derived exchange term is
\begin{align}
 B_{\mathrm{ex}}(y;P)
 &=s_{\mathrm{ex}}\sum_i
       \bigl[y_i(1-a_{ia})+(1-y_i)(1-a_{ib})\bigr] \nonumber\\
 &\quad+s_{\mathrm{ex}}\sum_{i<j}(1-2p_{ij})
       (1-y_i-y_j+2y_iy_j).
 \label{eq:matrix-exchange}
\end{align}
The first term favors assigning a customer to the route whose fixed members
have higher predicted affinity. The second rewards a shared assignment when
$p_{ij}>1/2$ and penalizes it when $p_{ij}<1/2$. Expanding these expressions
gives linear and quadratic QUBO coefficients directly from $P$.
In particular, let $w_{ij}=s_{\mathrm{ex}}(1-2p_{ij})$, extended
symmetrically. The bias adds
\begin{equation}
 \Delta h_i=\alpha\left[s_{\mathrm{ex}}(a_{ib}-a_{ia})
                   -\sum_{j\ne i}w_{ij}\right],
 \qquad \Delta J_{ij}=2\alpha w_{ij}
 \label{eq:bias-coefficients}
\end{equation}
to an upper-triangular QUBO
$c+\sum_i h_i y_i+\sum_{i<j}J_{ij}y_iy_j$. The remaining contribution
is constant. The mapping requires no extra decision variables and no
neural evaluation inside the solver. For example, $p_{ij}=0.8$ gives
a negative shared-route coefficient, while $p_{ij}=0.2$ gives a positive
coefficient of the same magnitude.

For reorders, define $\phi(u,v)=1-p_{uv}$ when both endpoints are customers
and $\phi(u,v)=0$ when either is the depot. The corresponding term is
\begin{align}
 B_{\mathrm{re}}(x;P)=s_{\mathrm{re}}\Bigl[
 &\sum_i\phi(s,i)x_{i1}
 +\sum_{\ell=1}^{k-1}\sum_{i\ne j}\phi(i,j)x_{i\ell}x_{j,\ell+1}
 \nonumber\\
 &+\sum_i\phi(i,e)x_{ik}\Bigr].
 \label{eq:matrix-reorder}
\end{align}
The scales $s_{\mathrm{ex}}$ and $s_{\mathrm{re}}$ are the means of positive
insertion costs and positive local distances, respectively, with a unit
fallback if none are positive. This uses same-route predictions as a soft
preference for consecutive visits; it does not assume that the matrix
predicts route adjacency.

The implementation also supports confidence mode, multiplying each
probability-derived cost by $|2p-1|$ for its corresponding probability
(including route affinities). Penalties in $\Pi_\alpha$ are increased to
cover the additional bias, preserving the underlying constraint treatment.
Setting $\alpha=0$ recovers the geometric QUBO control. Thus, comparing
positive $\alpha$ with zero measures the effect of the proposed matrix
coupling without changing the basic decision variables. Raw energies across
these objectives are not directly comparable; true route cost and feasibility
remain the evaluation criteria.

\subsection{Refinement solvers and classical controls}

Classical controls use 2-opt for reorders and relocate/swap for exchanges,
optionally with multiple starting points. Exact controls use Held--Karp
dynamic programming and assignment enumeration. Simulated annealing (SA)
is a classical QUBO sampler. Simulated quantum annealing (SQA) uses
path-integral Monte Carlo with 16 Trotter replicas and also runs on a classical
computer.

QAOA is evaluated using an exact, noise-free statevector simulator.
The cost unitary encodes the QUBO energy and alternates with an $X$ mixer.
For $v$ binary variables, substituting $z_i=(1-Z_i)/2$ gives a diagonal
cost Hamiltonian $H_C$. With $H_X=\sum_i X_i$, the depth-$p$ state is
\begin{equation}
 |\gamma,\beta\rangle=
 e^{-\mathrm i\beta_p H_X}e^{-\mathrm i\gamma_p H_C}\cdots
 e^{-\mathrm i\beta_1 H_X}e^{-\mathrm i\gamma_1 H_C}
 |+\rangle^{\otimes v}.
 \label{eq:qaoa-state}
\end{equation}
The final loop uses depth one, COBYLA angle optimization from four random
starts, 1,024 shots, and at most 16 qubits. Up to 50 lowest-energy distinct
sampled states are decoded and scored on true route cost. The random control
uses uniform bitstrings with the same shot count, retention rule, repair, and
qubit limit. Larger QUBOs are skipped unchanged. The shot budget concerns
final sampling; angle optimization uses exact expectations rather than noisy
hardware measurements.

\section{Experimental protocol}
\label{sec:protocol}

\subsection{Data and evaluation panels}

The standard synthetic recipe samples depot and customer coordinates
uniformly in $[0,1]^2$ and integer demands uniformly from
$\{1,\ldots,9\}$. The generator also supplies clustered (C) and mixed
random/clustered (RC) coordinates, alongside random (R) coordinates, at
the same three customer counts. Training uses the audited per-size pools
described in Section~\ref{sec:problem}. The spatial stress panel selects
eight stable-label graphs from each regime--size cell; it therefore
measures robustness on those selected cells rather than the full unfiltered
spatial distribution.


Development uses 72 graphs, 24 per size. A separate validation set contains
45 graphs, with 16, 16, and 13 at sizes 20, 50, and 100. These sets informed
refinement settings and exploratory ablations. The final refinement protocol
was declared before refinement runs on 96 reserved graphs, 32 per size, and
72 spatial stress graphs, eight per regime and size. The starting baseline
had already been scored and frozen. Each final configuration was then run
once with seed zero, all four selectors, and a maximum of 30 rounds.

The reserved reference is the best of the original label and two seeded
PyVRP runs with 10, 20, and 40 seconds of search per size, respectively.
Other panels use their instance labels. PyVRP is a strong hybrid genetic
search reference solver~\cite{wouda2024pyvrp}.

\subsection{Metrics and uncertainty}

Reported gaps are mean
instance-relative gaps:
\begin{equation}
 g_i=100\frac{C_i-C_i^{\mathrm{ref}}}{C_i^{\mathrm{ref}}},\qquad
 \bar g=\frac{1}{N}\sum_{i=1}^N g_i.
\end{equation}
They are not percentages above a proven optimum. Paired comparisons report
the mean of $g_i^{(a)}-g_i^{(b)}$ in percentage points (pp), with 95\%
percentile-bootstrap intervals from 10,000 graph resamples, seed zero.
A negative difference favors the first method. Intervals describe graph
variation under fixed trained models and do not measure training-seed
uncertainty or adjust for multiple comparisons.

For fixed subproblems, the local reference is the exact best reorder or
the best feasible assignment under the shared exchange reconstruction
rule. A hit means reaching that local reference cost. Kept excess is
$100(C_{\mathrm{kept}}-C_{\mathrm{local}}^*)/C_{\mathrm{local}}^*$,
where the original incumbent is retained if no candidate improves it.
Local excess, full-instance reference gap, and route-cost reduction have
different denominators and are reported separately.

\subsection{Refinement budgets and controls}

The notation $k/m/h$ denotes maximum reorder length, movable exchange
customers, and neighborhoods per selector per round. ``10 starts'' means
the current solution plus nine random starts, matching the code's
\texttt{classical\_restarts=10}. ``One start'' is a single local-search run
per subproblem, repeated across refinement rounds, not just one loop round.
SA uses 32 reads of 100 sweeps in the final loops. The matrix-derived
objective terms are evaluated in the development and validation bias
ablations. The final SA and QAOA arms use $\alpha=0$, while retaining
matrix-guided neighborhood selection. Accordingly, the final-test tables
establish comparisons between refinement operators; they do not constitute
a reserved-test evaluation of the positive-$\alpha$ formulation.

The fixed-subproblem comparisons use identical extracted subproblems. Full
loops share starting routes and neighborhood rules, but accepted moves cause
their later neighborhoods to diverge. Equal neighborhood limits therefore do
not imply identical search trajectories or equal wall-clock budgets.

\section{Constructive solver results}
\label{sec:construction-results}

Table~\ref{tab:construction} reports the frozen neural architecture before
refinement. Every output is feasible. The reserved panel has the strongest
references, making it the primary construction result: mean gaps are
8.56\%, 21.88\%, and 27.20\% at the three sizes, and 19.22\% overall.
The growth with customer count shows that feasibility and efficient
construction alone do not guarantee high route quality.

\begin{table}[tbp]
\centering\small
\caption{Mean reference gap (\%) of the frozen constructive solver.
Development is used for selection; validation also informs refinement
settings. Reserved references are strengthened PyVRP solutions, while
other panels use their instance labels.}
\label{tab:construction}
\begin{tabular}{@{}lrrrrr@{}}
\toprule
Panel & Graphs & N20 & N50 & N100 & All\\
\midrule
Development & 72 & 9.19 & 21.51 & 24.73 & 18.48\\
Validation & 45 & 8.31 & 21.08 & 27.26 & 18.33\\
Reserved & 96 & 8.56 & 21.88 & 27.20 & 19.22\\
Spatial R/C/RC & 72 & 8.29 & 19.20 & 22.09 & 16.53\\
\bottomrule
\end{tabular}
\end{table}

The constructive policy is materially better than directly decoding the
diffusion matrix into routes. On the same 72 development graphs and
cached prior, diffusion-only decoding gives gaps of 21.33\%, 29.61\%,
and 28.51\%, compared with the policy's 9.19\%, 21.51\%, and 24.73\%.
The paired policy-minus-diffusion differences are approximately
$-12.1$ pp $[-17.6,-6.6]$, $-8.1$ pp $[-12.9,-3.4]$, and
$-3.8$ pp $[-7.2,-0.5]$. This supports the policy's value on that panel;
it does not separately identify the effects of fusion, the local encoder,
or multi-start decoding.

The policy also improves over diffusion-only decoding in all nine spatial
regime--size cells. Its clustered-cell gaps are 10.51\%, 15.30\%, and
16.08\%; the corresponding random-cell gaps are 8.23\%, 22.54\%, and
26.67\%. Thus clustered instances are easier at N50 and N100, while N20
does not follow that ordering. The selected spatial panel shows that the
implementation works across the tested coordinate regimes, with the
selection caveat described in Section~\ref{sec:protocol}.

The learned solver remains behind strong classical construction-and-search.
On the same development panel, one-second PyVRP runs reach gaps of
0.00\%, 0.09\%, and 1.11\%. This comparison establishes practical headroom;
it is not a matched-hardware claim that the neural architecture is faster
or better than PyVRP. The following experiments test how much of that
headroom the local refinement interface can recover.

\section{Final refinement results}

\begin{table}[tbp]
\centering\small
\setlength{\tabcolsep}{3.5pt}
\caption{Final mean reference gaps (\%). The reserved set has 96 graphs;
the R/C/RC set has 72. All configurations use the same frozen starting
solutions within each set. Bold values identify the lowest gaps.}
\label{tab:final}
\begin{tabular}{@{}llcrrrrr@{}}
\toprule
 & & & \multicolumn{4}{c}{Reserved test} & R/C/RC\\
ID & Refinement & $k/m/h$ & N20 & N50 & N100 & All & All\\
\midrule
-- & Baseline & -- & 8.56 & 21.88 & 27.20 & 19.22 & 16.53\\
T1 & Classical, one start & 5/8/10 & 3.85 & 9.59 & 15.63 & 9.69 & 7.84\\
T2 & Classical, 10 starts & 5/8/30 & 1.92 & 5.81 & 8.89 & 5.54 & 5.00\\
T3 & Classical, 10 starts & 10/16/30 & \textbf{1.16} & \textbf{4.48} & \textbf{6.77} & \textbf{4.14} & \textbf{3.54}\\
T4 & SA & 5/8/30 & 3.05 & 6.36 & 9.36 & 6.26 & 5.23\\
T5 & SA + classical & 5/8/30 & 1.89 & 5.18 & 8.19 & 5.09 & 4.49\\
T6 & QAOA, depth 1 & 4/6/10 & 2.36 & 6.77 & 12.24 & 7.12 & 6.20\\
T7 & Random QUBO & 4/6/10 & 2.43 & 7.57 & 12.36 & 7.46 & 6.15\\
T8 & Classical, 10 starts & 4/6/10 & 2.52 & 6.82 & 12.24 & 7.19 & 6.34\\
T9 & QAOA + classical & 4/6/10 & 2.24 & 6.67 & 12.05 & 6.99 & 6.06\\
T10 & Classical exact & 4/6/10 & 2.21 & 6.74 & 12.13 & 7.03 & 6.10\\
\bottomrule
\end{tabular}
\end{table}

Table~\ref{tab:final} shows that local refinement improves the learned
solutions substantially. T1 reduces the reserved mean gap from 19.22\%
to 9.69\%, while T3 reaches 4.14\%. The analogous spatial stress gaps fall
from 16.53\% to 7.84\% and 3.54\%. Every final configuration improves
95 of 96 reserved graphs and all 72 spatial stress graphs. The remaining
reserved graph already matched its reference. Accepted outputs stay feasible.

The strongest observed configuration is classical search with larger
neighborhoods (T3). This is a best-configuration comparison, not evidence
isolating the solver: T3 is allowed larger moves than the QAOA arms. The
matched-limit comparisons are T4 versus T2 and T6 versus T7, T8, and T10.

\begin{table}[tbp]
\centering\small
\caption{Predeclared paired gap differences in pp, with 95\% bootstrap
intervals. Negative values favor the first method. An interval containing
zero is inconclusive, not a demonstration of equivalence.}
\label{tab:contrasts}
\begin{tabular}{@{}lrr@{}}
\toprule
Comparison & Reserved test & R/C/RC\\
\midrule
T1 $-$ baseline & $-9.53\ [-10.54,-8.49]$ & $-8.68\ [-9.90,-7.57]$\\
T3 $-$ baseline & $-15.08\ [-16.47,-13.65]$ & $-12.98\ [-14.54,-11.43]$\\
T4 $-$ T2: SA $-$ classical & $+0.72\ [+0.16,+1.27]$ & $+0.23\ [-0.33,+0.79]$\\
T5 $-$ T2: SA before polish & $-0.45\ [-0.96,+0.04]$ & $-0.51\ [-1.02,-0.01]$\\
T6 $-$ T7: QAOA $-$ random & $-0.34\ [-0.76,+0.02]$ & $+0.05\ [-0.10,+0.20]$\\
T6 $-$ T10: QAOA $-$ exact & $+0.09\ [-0.18,+0.35]$ & $+0.09\ [-0.06,+0.29]$\\
T6 $-$ T8: QAOA $-$ classical & $-0.07\ [-0.56,+0.38]$ & $-0.15\ [-0.53,+0.23]$\\
T9 $-$ T8: QAOA before polish & $-0.21\ [-0.68,+0.24]$ & $-0.29\ [-0.63,+0.03]$\\
\bottomrule
\end{tabular}
\end{table}

SA is worse than restarted classical search on the reserved test by
0.72 pp, with an interval entirely above zero (Table~\ref{tab:contrasts}).
On the spatial stress set this comparison is inconclusive. Adding SA before
classical polish reduces the mean gap by about half a point, but the reserved
interval includes zero. The spatial interval narrowly excludes zero; this
small result should be interpreted alongside the multiple comparisons and
the additional computation. In the archived final runs, the hybrid consumes
approximately 26--39\% more solver time than its classical control.

QAOA gives no statistically resolved improvement over random sampling,
restarted classical search, or exact classical subproblem solving on either
final set. Its reserved gap is 7.12\%, compared with 7.46\%, 7.19\%, and
7.03\%, respectively. QAOA followed by classical polish also has an interval
including zero against classical search alone. Approximately 0.6--0.8\% of
QAOA subproblem attempts exceed the 16-qubit limit and are skipped.

\section{Understanding the benchmark outcomes}

\subsection{Formulation and single-subproblem performance}

The reorder QUBO's exact minimum agrees with the true optimum on all 84
enumerable development and validation segments. For exchanges, exact QUBO
minimization returns a capacity-feasible assignment on all 698 tested
subproblems, but matches the best reconstructed route cost on only 197/430
development and 117/268 validation exchanges. About one fifth of those
surrogate-optimal candidates are worse than the starting routes. This
distinction limits any solver that optimizes the exchange QUBO, even if it
finds its ground state.

\begin{table}[tbp]
\centering\small
\caption{Local optimum hit rates (\%) on frozen subproblems. There are
392/243 development/validation reorders and 430/268 exchanges. Budgeted
methods target 0.25 s per subproblem; individual work units can overrun it.
The SQA comparison is a separate run, with an SA control from that run.}
\label{tab:local}
\begin{tabular}{@{}lrrrr@{}}
\toprule
 & \multicolumn{2}{c}{Reorder} & \multicolumn{2}{c}{Exchange}\\
Method & Dev. & Val. & Dev. & Val.\\
\midrule
Classical, one start & 92.9 & 95.9 & 76.5 & 76.9\\
Classical, budgeted restarts & 100.0 & 100.0 & 91.9 & 91.8\\
SA, budgeted & 50.0 & 53.9 & 67.7 & 69.0\\
SA + prior bias, $\alpha=0.5$ & 52.0 & 62.6 & 70.9 & 74.6\\
\midrule
SA, separate SQA comparison & 60.5 & 59.3 & 74.0 & 71.6\\
SQA, same comparison & 27.6 & 25.9 & 35.1 & 35.4\\
\bottomrule
\end{tabular}
\end{table}

Classical restarts dominate the main local comparisons
(Table~\ref{tab:local}). SA is effective on very small reorders but becomes
less reliable at five customers, where the one-hot encoding already needs
25 variables. SQA performs worse than its SA control at the same nominal
time budget; simulating 16 replicas allows only a few reads per subproblem.
Approximately half of SA exchange reads violate capacity before repair.
Feasible final solutions therefore do not imply that the sampler itself
usually finds feasible assignments.

\subsection{What QAOA adds on small subproblems}

\begin{table}[tbp]
\centering\small
\caption{Optimum hit rates (\%) after shared decoding, repair, and retention
of the better of candidate and start, for exchanges of at most 16 qubits.
These screening results are exploratory, separate from the final loop.}
\label{tab:qaoa}
\begin{tabular}{@{}lrr@{}}
\toprule
Method & Development ($N=192$) & Validation ($N=119$)\\
\midrule
Classical, one start & 89.1 & 93.3\\
Classical, budgeted restarts & 97.4 & 99.2\\
SA & 72.4 & 74.8\\
SA + prior bias & 81.2 & 87.4\\
QAOA, depth 1 & 87.0 & 86.6\\
QAOA, depth 2 & 83.9 & 82.4\\
QAOA, depth 3 & 83.9 & 82.4\\
Uniform random sampling & 90.6 & 95.8\\
\bottomrule
\end{tabular}
\end{table}

QAOA can increase the probability of low-energy states without improving
the final routing result. On all 55 development and 29 validation reorders
small enough for screening, QAOA and matched random sampling both reach
100\% optimum hits after the shared decoding procedure. On exchanges,
random sampling has higher aggregate hit rates than the tested QAOA depths
(Table~\ref{tab:qaoa}). For tiny state spaces, 1,024 samples and strong
classical repair leave little room for a quantum sampling benefit.

Deeper cold-start circuits do not help in this screening. A separate
development warm-start experiment improves QUBO-optimal sampling by
initializing deeper circuits from shallower solutions. That diagnostic was
not used in the final depth-one loop and does not establish an improvement
in end-to-end routing. The exchange surrogate remains a separate obstacle.

\subsection{Evidence for the matrix-based formulation and selection}

The uncertainty selector produces 1.4--2.0 pp more first-round route-cost
reduction than nearby random neighborhoods matched in number and size,
across development and validation with both classical search and SA. Its
signal is useful for ranking possible customer moves. The low-confidence
edge selector does not show the same benefit: it identifies possible
membership errors but feeds reorder moves that cannot change membership.

The bias ablation directly evaluates the matrix-to-QUBO contribution in
Section~\ref{sec:matrix-qubo}. In a three-seed exploratory sweep,
probability-mode bias with $\alpha=2$ lowers kept excess over the local
reference by 0.84 pp on development and 1.27 pp on validation. Reorders
show no consistent benefit. A subsequent full-loop test reaches 5.72\%
and 6.51\% gap with biased SA, against 5.81\% and 6.20\% for classical
restarts at the same default limits. Those comparisons do not distinguish
the methods. The bias strength was chosen from the sweep, so this remains
development evidence rather than a final-test claim.

Increasing the number of rounds eventually saturates: all development and
validation extension runs stop within 15 rounds. Larger neighborhoods help
classical search more than SA. At limits 10/16/30, restarted classical
search reaches 4.23\%/4.10\% development/validation gap, compared with
6.84\%/6.77\% for SA. This motivates the larger classical arm T3.

\section{Computational cost and discussion}
\subsection{Measured runtime}

\begin{table}[tbp]
\centering\small
\caption{Mean refinement wall-clock seconds per graph in the separate
timing pass, four development graphs per size. Each run uses one process
and one numerical-library thread, with at most four runs concurrent.
These are not the final-test quality panel or quantum hardware timings.}
\label{tab:timing}
\begin{tabular}{@{}llrrr@{}}
\toprule
ID & Configuration & N20 & N50 & N100\\
\midrule
T1 & Classical, one start & 0.2 & 1.0 & 2.9\\
T2 & Classical, 10 starts & 1.5 & 15.6 & 49.3\\
T3 & Classical, larger neighborhoods & 3.3 & 26.2 & 134.6\\
T4 & SA & 3.8 & 17.9 & 47.2\\
T5 & SA + classical & 5.1 & 22.2 & 79.1\\
T6 & Simulated QAOA & 40.1 & 159.2 & 145.7\\
T7 & Random QUBO & 0.9 & 5.1 & 8.2\\
T10 & Classical exact & 0.5 & 3.0 & 5.0\\
\bottomrule
\end{tabular}
\end{table}

The economical classical option takes approximately 0.2--3 seconds per
graph; larger restarted search takes about 3 seconds to 2.2 minutes
(Table~\ref{tab:timing}). Add approximately 0.78, 0.86, and 1.06 seconds
for the frozen prior and policy at the three sizes on an RTX 3060 Ti,
excluding model loading. Most archived refinement runs used ten workers
sharing six physical CPU cores; their solver times should not be treated
as isolated deployment timings. In those final reserved runs, QAOA takes
244.6 s per graph versus 5.9 s for exact classical solving, approximately
41 times as long. This ratio describes classical simulation only.

\subsection{What the architecture learns and what search repairs}

The empirical picture separates construction from improvement. The policy
uses the diffusion prior more effectively than direct matrix decoding, but
its initial solutions still contain substantial route-cost headroom.
Exchange moves correct membership errors and can reduce the number of
routes; reorder moves improve visit order without changing membership.
In exploratory three-round development loops, N50 starts with about 7.9
routes per instance and ends with roughly 6.9 under the stronger settings.
This is consistent with refinement repairing route fragmentation, although
it does not isolate a single architectural cause of that fragmentation.

The prior has three distinct interfaces: local policy attention,
neighborhood ranking, and optional QUBO objective bias. Evidence at one
interface cannot be transferred automatically to the others. Useful
exchange ranking does not prove that hard attention masks are optimal,
and improvements from positive matrix bias on development do not establish
an improvement on reserved graphs. Conversely, setting objective bias to
zero in the final tests does not remove matrix guidance from the whole
pipeline: the starting policy and two selectors still depend on $P$.

The strongest observed refinement is classical search with larger
neighborhoods. This points to a practical configuration for the current
implementation: use the frozen constructive policy to obtain a feasible
incumbent, then choose the classical neighborhood budget according to the
available latency. It also identifies a research direction: improve the
constructive policy and the exchange representation before treating a
more sophisticated QUBO sampler as the main bottleneck.

\subsection{Limits of the evidence}

The study evaluates selected checkpointed models, not all possible versions
of the architecture. In particular, it does not present an independent
reserved-panel factorial ablation of global-only attention, local attention,
gated fusion, savings bias, and training-start count. The observed pipeline
improvements should therefore not be attributed to the gate or any other
single component without a matched retraining experiment.

The spatial panel uses stable-label selections from larger pools and had
informed earlier baseline diagnostics. It is not a wholly untouched
population-level generalization test. Reference costs are heuristic, most
final configurations use one refinement seed, and graph bootstrap intervals
do not capture training variability. Equal neighborhood limits also do not
establish equal compute: accepted moves lead to different later subproblems,
and fixed-read and nominal-time-budget comparisons measure different
resources.

For the QUBO extension, the key limitations are the small simulated state
spaces, the surrogate exchange objective, and the contribution of classical
repair. A better QUBO-energy distribution can still produce worse route
cost after reconstruction. Statevector QAOA ignores hardware noise,
connectivity, compilation, and finite-shot parameter-estimation costs.
SQA also executes entirely on a classical computer. The measured timings
therefore do not predict deployment on a quantum device.

The prior-guidance ablations do not determine whether the diffusion stage
justifies its total training and inference cost relative to a fully
classical construction-and-search pipeline. Large-scale instances,
fixed-fleet constraints, and time windows are outside the reported
architecture evaluation. Future extensions should first establish stronger
construction baselines and independent component ablations, then evaluate
positive-bias refinement with a frozen protocol and an explicitly matched
compute budget.

\section{Reproducibility and artifact provenance}
\label{sec:reproducibility}

The constructive solver is frozen as \texttt{baseline-v1.0}. Its manifest
pairs source commit \texttt{460b2efb5} with SHA-256 hashes for all six
checkpoints, copied training configurations, dataset identities, and
panel manifests. The frozen evaluation entry point was committed in
\texttt{aca6c0aec}. Prior inference uses a graph-specific seed derived
from the first four bytes of SHA-256 of the base seed and instance ID,
with base seed zero. This avoids making a graph's prior depend on its
position in a batch.

The final refinement protocol was declared on 2 October 2026 before
refinement runs on the reserved and spatial panels. It fixes T1--T10,
selector limits, seed zero, stopping rules, and paired comparisons.
Evidence archives record aggregate results and analysis scripts; full
per-graph rows and large generated checkpoints remain in local output
directories. The article itself is self-contained and requires neither
external figures nor a separate bibliography file to compile.

The archived training environment records Python 3.12.10, PyTorch
2.11.0+cu128, CUDA 12.8, and an NVIDIA GeForce RTX 3060 Ti with 8 GiB of
memory on Windows 11. Refinement reporting records six physical CPU cores
and the concurrency described above, but the retained evidence does not
provide a complete pinned software environment and CPU model for every
benchmark pass. Project dependency constraints are specified in
\artifact{pyproject.toml}; those constraints should not be confused with
an exact historical lockfile. Appendix~\ref{app:artifacts} maps the
reported methods and results to retained project artifacts.

\section{Conclusion}

We have presented our implemented diffusion-guided CVRP solver, inspired
by the structural-prior idea of CMD. A discrete route-membership denoiser
feeds a global and local Transformer policy, learned gated fusion, and
a capacity-aware dual-pointer decoder. Multi-start inference produces
feasible solutions, while a modular local improvement stage reuses the
prior to select exchanges and to form optional matrix-informed QUBOs.

The complete system is functional across the tested 20-, 50-, and
100-customer panels. Its constructive reserved gap is 19.22\%; inexpensive
classical refinement reduces it to 9.69\%, and larger restarted search
to 4.14\%. The corresponding spatial stress gap falls from 16.53\% to
3.54\%. These outcomes establish substantial refinement value while
showing that the current learned construction remains behind strong
classical CVRP solving.

Matrix-guided selection identifies useful exchanges, and objective bias
improves annealing in exploratory local comparisons. Simulated QAOA gives
no reliable benefit over matched random sampling or exact classical
solving at its small neighborhood limits. The architecture's practical
contribution is a complete, inspectable combination of structural
prediction, feasible construction, and solver-independent refinement.
Further claims about individual components, positive-bias generalization,
or quantum advantage require additional controlled evidence.

\begin{thebibliography}{99}

\bibitem{wang2026cmd}
Z. Wang, T. Zhang, N. Xue, E. Ozcan, L. Wang, and R. Bai.
\newblock Constraints Matrix Diffusion based Generative Neural Solver for
Vehicle Routing Problems.
\newblock arXiv:2603.07568v1, 2026.
\newblock \url{https://arxiv.org/abs/2603.07568v1}.

\bibitem{kool2019attention}
W. Kool, H. van Hoof, and M. Welling.
\newblock Attention, Learn to Solve Routing Problems!
\newblock \emph{International Conference on Learning Representations}, 2019.
\newblock \url{https://arxiv.org/abs/1803.08475}.

\bibitem{kwon2020pomo}
Y.-D. Kwon, J. Choo, B. Kim, I. Yoon, Y. Gwon, and S. Min.
\newblock POMO: Policy Optimization with Multiple Optima for Reinforcement
Learning.
\newblock \emph{Advances in Neural Information Processing Systems}, 2020.
\newblock \url{https://arxiv.org/abs/2010.16011}.

\bibitem{austin2021discrete}
J. Austin, D. D. Johnson, J. Ho, D. Tarlow, and R. van den Berg.
\newblock Structured Denoising Diffusion Models in Discrete State-Spaces.
\newblock \emph{Advances in Neural Information Processing Systems}, 2021.
\newblock \url{https://arxiv.org/abs/2107.03006}.

\bibitem{sun2023difusco}
Z. Sun and Y. Yang.
\newblock DIFUSCO: Graph-based Diffusion Solvers for Combinatorial Optimization.
\newblock \emph{Advances in Neural Information Processing Systems}, 2023.
\newblock \url{https://arxiv.org/abs/2302.08224}.

\bibitem{vidal2022hgs}
T. Vidal.
\newblock Hybrid Genetic Search for the CVRP: Open-Source Implementation and
SWAP* Neighborhood.
\newblock \emph{Computers \& Operations Research}, 140:105643, 2022.
\newblock \url{https://doi.org/10.1016/j.cor.2021.105643}.

\bibitem{wouda2024pyvrp}
N. A. Wouda, L. Lan, and W. Kool.
\newblock PyVRP: A High-Performance VRP Solver Package.
\newblock \emph{INFORMS Journal on Computing}, 36(4):943--955, 2024.
\newblock \url{https://doi.org/10.1287/ijoc.2023.0055}.

\bibitem{held1962dynamic}
M. Held and R. M. Karp.
\newblock A Dynamic Programming Approach to Sequencing Problems.
\newblock \emph{Journal of the Society for Industrial and Applied
Mathematics}, 10(1):196--210, 1962.
\newblock \url{https://doi.org/10.1137/0110015}.

\bibitem{lucas2014ising}
A. Lucas.
\newblock Ising Formulations of Many NP Problems.
\newblock \emph{Frontiers in Physics}, 2:5, 2014.
\newblock \url{https://doi.org/10.3389/fphy.2014.00005}.

\bibitem{farhi2014qaoa}
E. Farhi, J. Goldstone, and S. Gutmann.
\newblock A Quantum Approximate Optimization Algorithm.
\newblock arXiv:1411.4028, 2014.
\newblock \url{https://arxiv.org/abs/1411.4028}.

\bibitem{azad2020vrp}
U. Azad, B. K. Behera, E. A. Ahmed, P. K. Panigrahi, and A. Farouk.
\newblock Solving Vehicle Routing Problem Using Quantum Approximate
Optimization Algorithm.
\newblock \emph{IEEE Transactions on Intelligent Transportation Systems}, 2022.
\newblock \url{https://doi.org/10.1109/TITS.2022.3172241}.

\end{thebibliography}

\appendix
\section{Mapping the article to project artifacts}
\label{app:artifacts}

All paths in Table~\ref{tab:artifacts} are relative to the repository root.
External references supply methodological context; the numerical routing
results come from the project archives listed here. No new experiments
are introduced by this article.

\small
\begin{longtable}{@{}p{0.30\linewidth}p{0.65\linewidth}@{}}
\caption{Implementation and evidence sources. Folder braces denote the
named alternatives within that folder.}\label{tab:artifacts}\\
\toprule
Article content & Artifact\\
\midrule
\endfirsthead
\toprule
Article content & Artifact\\
\midrule
\endhead
\bottomrule
\endfoot
Problem representation and data &
\artifact{src/vrp_diffusion_quantum/data/generate_cvrp.py};
\artifact{src/vrp_diffusion_quantum/data/training_labels.py}\\[4pt]
Diffusion schedule and denoiser &
\artifact{src/vrp_diffusion_quantum/models/diffusion.py};
\artifact{src/vrp_diffusion_quantum/models/constraint_denoiser.py}\\[4pt]
Global, local, and fusion modules &
\artifact{src/vrp_diffusion_quantum/models/}\newline
\artifact{global_encoder.py}, \artifact{local_masked_encoder.py},
\artifact{fusion_encoder.py}\\[4pt]
Decoding and training &
\artifact{src/vrp_diffusion_quantum/models/decoder.py};
\artifact{src/vrp_diffusion_quantum/train/}\newline
\artifact{train_diffusion.py}, \artifact{train_policy.py}\\[4pt]
Frozen architecture and construction results &
\artifact{docs/evidence/baseline_freeze_20261001/freeze_manifest.json};
\artifact{docs/evidence/baseline_freeze_20261001/configs/};
\artifact{docs/classical_baseline_freeze_proposal_2026-09-30.md}\\[4pt]
Final evaluation declaration &
\artifact{docs/phase6_final_test_protocol_2026-10-02.md}\\[4pt]
Final gaps and paired contrasts &
\artifact{docs/evidence/phase6_20261001/final/analysis_reserved.json};
\artifact{docs/evidence/phase6_20261001/final/analysis_ood.json};
corresponding run \artifact{summary.json} files\\[4pt]
QUBO coefficients, repair, and acceptance &
\artifact{src/vrp_diffusion_quantum/quantum/}\newline
\artifact{qubo_reorder.py}, \artifact{qubo_exchange.py},
\artifact{qubo_bias.py}, \artifact{refinement.py}\\[4pt]
Local solver comparisons &
\artifact{docs/evidence/phase6_20261001/}\newline
\artifact{experiment_{development,validation}/summary.json};
\artifact{experiment_{development,validation}_sqa/summary.json}\\[4pt]
QAOA and random screening &
\artifact{docs/evidence/phase6_20261001/}\newline
\artifact{qaoa_random_control_{development,validation}.json}\\[4pt]
Matrix guidance and larger-neighborhood studies &
\artifact{docs/evidence/phase6_20261001/}\newline
\artifact{bias_sweep/}, \artifact{selection/}, \artifact{extension/};
\artifact{docs/phase6_refinement_experiments_2026-10-01.md}\\[4pt]
Separate timing pass &
\artifact{docs/evidence/phase6_20261001/timing/}\newline
\artifact{{t1_t10,t2_t4_t5,t3,t6_t7,exact_small}/summary.json}\\
\end{longtable}
\normalsize

\end{document}
